\tzpanchor{0401}
\tzpentryheader{0401}{Electromagnetic Stress-Energy Tensor}

\begingroup\small
\begingroup\scriptsize
\setlength{\tabcolsep}{4pt}
\renewcommand{\arraystretch}{0.92}
\noindent\begin{tabularx}{\linewidth}{@{}p{0.24\linewidth}p{0.36\linewidth}X@{}}
\textbf{UUID} & \multicolumn{2}{l@{}}{\tzpcode{729ad140-86a2-5b8c-8960-62ffe99a1ed8}}\\
\textbf{PiID} & \tzpcode{3057270365759} & \textbf{TZPi}\quad \tzpcode{0}\quad \textbf{PiPE}\quad \tzpcode{408}\\
\textbf{RTEID} & \textbf{Poloidal $\theta$}\quad \tzpcode{1557.1721 rad} & \textbf{Toroidal $\phi$}\quad \tzpcode{02.5657 rad}\\
\textbf{TZPID} & \tzpcode{ID0401} & \textbf{Node Dependencies}\quad \tzpidref{0400}\\
\textbf{Registry} & \tzpcode{registry\_id\_0401} & \textbf{Scale}\quad Quantum-Classical Transition\\
\textbf{LOC Classification} & QA641 manifolds and topology & QC174.17 mathematical physics\\
\textbf{arXiv Classification} & Primary: math-ph & Cross-list: gr-qc, math.DG\\
\textbf{Primary Support} & Aggregate Relation & Support Relation\\
\textbf{Closure Support} & Closure Relation & \\
\end{tabularx}\par
\endgroup
\endgroup

\tzpsection{Canonical Statement}
ID 401 maps how the vacuum pressure residue is converted into momentum flux.

\tzpsection{Canonical Equation}
\[
\begin{adjustbox}{max width=\linewidth,center}
$\displaystyle
\mathcal{L}_{\mathrm{bundle}}=\int_{\Sigma}F/(2\pi)
$
\end{adjustbox}
\]

\begin{minipage}[t]{0.49\linewidth}
\tzpsection{Canonical Symbol Key}
\begingroup
\small
\raggedright
\textbf{Source equation symbols.}\par
\begin{itemize}[leftmargin=1.35em,itemsep=1pt,topsep=1pt,parsep=0pt]
    \item $\mathcal{L}_{\mathrm{bundle}}$: source equation symbol.
    \item $F$: source equation symbol.
    \item $\pi$: source equation symbol.
    \item $x$: local field point or coordinate variable in the canonical equation.
    \item $S$: source equation symbol.
\end{itemize}
\endgroup
\end{minipage}\hfill
\begin{minipage}[t]{0.47\linewidth}
\tzpsection{Wolfram Form}
\begingroup
\scriptsize
\raggedright
\textbf{Source Wolfram model.}\par
\begin{Verbatim}[fontsize=\tiny,breaklines=true,breakanywhere=true]
(* TZPID ID0401 generated source Wolfram model *)
ClearAll[K0401, Cr0401, T, R, A, W, I, N];
eq0401$1 = HoldForm[$LBundle == (IntegralSigmaF/(2*Pi))*$];

K0401 = HoldForm[{eq0401$1}];

N = "normalized TRAWIN closure object for Electromagnetic Stress-Energy Tensor";
I = "Electromagnetic Stress-Energy Tensor integration/comparison layer";
W = "Electromagnetic Stress-Energy Tensor wave or operator carrier";
A = "Electromagnetic Stress-Energy Tensor admissible action/amplitude condition";
R = "Electromagnetic Stress-Energy Tensor rotation/curl preservation layer";
T = "Electromagnetic Stress-Energy Tensor local source condition";

Cr0401[expr_] := HoldForm[N[I[W[A[R[T[expr]]]]]]];

Result0401 = Cr0401[K0401];
\end{Verbatim}
\endgroup
\end{minipage}

\par\vspace{0.45\baselineskip}
\noindent\begin{minipage}[t]{\linewidth}
\begingroup
\small
\raggedright
\textbf{TRAWIN Operator Key.}\par
\noindent\textbf{Time/localization operator:} $\mathsf{T}\!\left[\mathrm{local}\right]$ Electromagnetic Stress-Energy Tensor local source condition.\par
\noindent\textbf{Rotation/curl operator:} $\mathsf{R}\!\left[\nabla\times\right]$ Electromagnetic Stress-Energy Tensor rotation/curl preservation layer.\par
\noindent\textbf{Action/amplitude operator:} $\mathsf{A}\!\left[\mathrm{admissible}\right]$ Electromagnetic Stress-Energy Tensor admissible action/amplitude condition.\par
\noindent\textbf{Wave/d'Alembertian operator:} $\mathsf{W}\!\left[\Box\right]$ Electromagnetic Stress-Energy Tensor wave or operator carrier.\par
\noindent\textbf{Integration operator:} $\mathsf{I}\!\left[\int\right]$ Electromagnetic Stress-Energy Tensor integration/comparison layer.\par
\noindent\textbf{Normalization operator:} $\mathsf{N}\!\left[\mathrm{normalized}\right]$ normalized TRAWIN closure object for Electromagnetic Stress-Energy Tensor.\par
\endgroup
\end{minipage}

\clearpage
\noindent\textbf{[ID 0401] Electromagnetic Stress-Energy Tensor}\par
\vspace{0.35em}
\tzpsection{Derivation Layer}
\paragraph{Primary Equation.}
\begin{equation}
\det(g)\big|_{p}=0 .
\end{equation}

\paragraph{Primary Derivation.}
Let $(M,g)$ be a manifold equipped with metric tensor $g$. Away from the Trawin Zero Point $p$, the metric is assumed nondegenerate:

\begin{equation}
\det(g)\neq 0,
\qquad x\neq p .
\end{equation}

At the TZP, the manifold undergoes a structural collapse of ordinary metric invertibility. This means the inverse metric $g^{-1}$ cannot be defined at $p$.

A metric tensor is invertible if and only if its determinant is nonzero. Therefore, non-invertibility at $p$ implies:

\begin{equation}
\det(g)\big|_{p}=0 .
\end{equation}

\paragraph{Interpretation.}
The TZP is marked by metric degeneracy. This is the first mathematical signature that ordinary geometry fails at the zero point and must be replaced by a Trawinistic structure.

\par\vspace{0.35em}
\noindent\textbf{Derivable Consequences.}\quad
\begingroup
\footnotesize
\raggedright
The canonical equation anchors Electromagnetic Stress-Energy Tensor by connecting the source condition to the derivation layer and proof certificate.
\par\endgroup

\tzpsection{Equation Support}
\noindent\textbf{Canonical Support Relation}\par
\[
\begin{adjustbox}{max width=\linewidth,center}
$\displaystyle
\mathcal{L}_{\mathrm{bundle}}=\int_{\Sigma}F/(2\pi)
$
\end{adjustbox}
\]
\noindent The support equation repeats the canonical relation so the derivation page remains self-contained.\par

\clearpage
\noindent\textbf{[ID 0401] Electromagnetic Stress-Energy Tensor}\par
\vspace{0.35em}
\tzpsection{Dictionary}
\begingroup\small
\tzpsection{Definition}
Electromagnetic Stress-Energy Tensor is a registry definition in the active TZPID arc.

\tzpsection{Contextual Meaning}
ID 401 maps how the vacuum pressure residue is converted into momentum flux.

\tzpsection{Formal Statement}
\[
\begin{adjustbox}{max width=\linewidth,center}
$\displaystyle
\mathcal{L}_{\mathrm{bundle}}=\int_{\Sigma}F/(2\pi)
$
\end{adjustbox}
\]

\tzpsection{Technical Interpretation}
ID 401 maps how the vacuum pressure residue is converted into momentum flux.

\tzpsection{Equation Grounding}
The equation grounding is carried by the canonical equation and the source Wolfram model on the entry page.

\tzpsection{Interpretive Role}
The canonical equation anchors Electromagnetic Stress-Energy Tensor by connecting the source condition to the derivation layer and proof certificate.

\tzpsection{TZP Thesaurus}
Electromagnetic Stress-Energy Tensor; canonical operator; source equation; TRAWIN-normalized relation.
\endgroup

\tzpsection{Encyclopedia Mathematica}
\begingroup\footnotesize\sloppy
The visual plate below is reserved for the source-specific equation and progression graphic for [ID 0401]. It is included only as a companion to the canonical equation, not as a substitute for the formal text.
\par\endgroup
\begingroup\footnotesize\itshape
Visual plate pending source-specific approval for [ID 0401].
\par\endgroup

\clearpage
\begingroup\tzpsupportsetup
\noindent\textbf{\fontsize{10.8pt}{12.8pt}\selectfont [ID 0401] Constructive Logic and Proof Certificate}\par
\noindent\textbf{\fontsize{10pt}{12pt}\selectfont Electromagnetic Stress-Energy Tensor}\par
\vspace{0.45em}
\begingroup
\footnotesize
\setlength{\parskip}{0.22em}
\setlength{\parindent}{0pt}
\noindent\textbf{Curry--Howard--Lambek Interpretation}\par
Under the Curry--Howard--Lambek correspondence, this registry entry is read as a typed proof
object: the stated source condition supplies the proposition, the derivation supplies the
morphism, and the proof certificate records the machine handles needed to check the obligation
in the formal registry.

\vspace{0.45em}
\noindent\textbf{CHL Proof}\par
\textbf{Statement:} t\textasciicircum{}\textbackslash{} = frac\textbackslash{}1\textbackslash{}\textbackslash{}\_0\textbackslash{} left[ F\textasciicircum{}\textbackslash{} F\textasciicircum{}\textbackslash{}\textbackslash{}\_\textbackslash{} - frac\textbackslash{}1\textbackslash{}\textbackslash{}4\textbackslash{} g\textasciicircum{}\textbackslash{} F\textbackslash{}\_\textbackslash{} F\textasciicircum{}\textbackslash{} right]

\textbf{Logic Validation:}\\
\textbf{Proposition:} ``ELECTROMAGNETIC STRESS-ENERGY TENSOR (T\textasciicircum{}\{\textbackslash{}mu\textbackslash{}nu) is valid as a tensorial/geometric
statement when the stated tensor fields are well formed on the specified domain.''\\
\textbf{Proof:} The source statement gives a metric, curvature, or tensor relation. The CHL reading
validates it by checking that the tensorial objects have compatible domains, indices,
and transformation behavior.

\textbf{VERDICT:} \ensuremath{\checkmark}\ \textbf{TRUE} --- tensorial well-formedness validation

\textbf{Formal Status:} \ensuremath{\checkmark}\ THEORETICAL -- source-backed CHL proof obligation

\textbf{Physical Status:} THEORETICAL -- requires model-specific empirical interpretation
\par\vspace{0.45em}
\noindent{\tiny\textbf{Isabelle/HOL.} \tzpplaincode{physics\_validity\_from\_model, external\_validity\_package, registry\_number\_matches, equation\_count\_matches}}\par
\vfill
\begingroup\footnotesize\itshape
Proof visual pending source-specific approval for [ID 0401].
\par\endgroup
\vfill
\endgroup
\endgroup
